\section{Introduction}
\label{sec:intro}
Multimodal models answer questions that cannot be resolved from any single input by combining images with text, audio,
or other signals~\cite{baltrusaitis2019multimodal,liang2024foundations}. Vision-language models (VLMs) now perform such
reasoning at scale~\cite{radford2021clip,alayrac2022flamingo,liu2023llava,bai2025qwen25vl}. A basic question for these
tasks is how much their answers require fusion: can the answer be determined from each modality separately, or does it
depend on how the modalities interact? Partial Information Decomposition (PID)~\cite{williams2010nonnegative,bertschinger2014quantifying}
is the standard framework for this question. It splits the information two modalities carry about a label into
redundant, unique and synergistic parts, where synergy is information available only from their combination. Synergy
is used to characterise multimodal datasets and guide model selection~\cite{liang2023quantifying,liang2024multimodal},
to analyse how VLMs use their inputs~\cite{xiu2026comprehensive,fang2026towards}, and to define training objectives
and curricula~\cite{kontras2026synib,yarici2026synco,singh2026spice}.

Synergy, however, is not the same as requiring fusion. Consider a task whose answer is ``yes'' only when both
modalities show an object (logical AND). Under a balanced binary setup, PID assigns this task 0.5 bits of synergy,
yet a model that scores each modality separately and adds the two scores answers it perfectly. What an additive model
cannot represent is a reversal: an XOR-like pattern in which the effect of one modality flips with the other. PID
cannot separate the two cases (Fig.~\ref{fig:teaser}). The reason is in its definition: BROJA PID~\cite{bertschinger2014quantifying}
fixes the relation between each modality and the label and optimises over the rest, so for deterministic labels every
PID term depends only on per-modality statistics, not on how answers are arranged across input pairs. That synergy and
interaction differ has been noted in information theory~\cite{james2019unique,jansma2023higher}, and model probes such
as EMAP~\cite{hessel2020does} show that many benchmarks can be solved additively. What multimodal evaluation lacks is a
model-free measure of how much a task requires non-additive fusion.

We bring such a measure from the classical theory of additive structure. For a task whose answer is known for every
pair of inputs, we arrange the answers in a table with one modality indexing rows and the other columns. The best
additive model is a \emph{rule}: one score per row and one per column, added and thresholded. The cells it gets wrong
are \emph{exceptions}. A classical result~\cite{ducamp1969composite,ryser1957combinatorial} says a table has no
exceptions exactly when it has no XOR square. We compute the rule exactly by integer programming, so the share of
exceptions is the minimum error of any additive model. It also sets a floor on interaction dependence: removing a
model's cross-modal interactions~\cite{hessel2020does} must cost at least the exception share minus the model's error.

We use this measure in two ways. First, on four VLMs~\cite{bai2025qwen25vl,wang2025internvl35,an2025llavaov15,laurencon2024building}
and 99 fully labelled answer tables from CLEVR, GQA, CLOVE and POPE~\cite{johnson2017clevr,hudson2019gqa,lei2023symbolic,li2023pope},
we ask how much each model relies on interactions relative to what its accuracy requires. Within each benchmark and at
fixed answer balance, the exception share ranks tables by interaction dependence and synergy adds little. An
accuracy-matched null scorer shows that most of this relation is the requirement itself: models depend on
interactions roughly as much as their accuracy demands, and one model (InternVL3.5) clearly more. Second, in continual
learning, forgetting is usually averaged over all samples~\cite{lopezpaz2017gradient,delange2021continual}; because
exceptions are a small fraction of the data, such averages cannot show whether a method keeps the knowledge that
requires fusion. On compositional CLEVR tasks, replay methods forget exceptions faster than rule cells with the same
answer and the same difficulty, and all methods learn new exceptions worse, while standard backward transfer shows
neither.

Our contributions are:
\begin{itemize}
  \item A model-free measure of how much a task requires non-additive fusion, computed exactly, with its link to
  additive models, to EMAP and to PID made explicit (Sec.~\ref{sec:theory}). The underlying characterisation is
  classical; the contribution is its use for multimodal evaluation.
  \item A test on four VLM families that separates the interaction dependence a model needs from the dependence it has,
  using structure--synergy decorrelated tables and an accuracy-matched null (Sec.~\ref{sec:vlm}).
  \item An exception-aware evaluation of twelve continual-learning methods with answer- and difficulty-matched
  comparisons, showing that on compositional tasks exceptions are forgotten faster and learned worse, and that standard
  backward transfer misses this (Sec.~\ref{sec:cl}).
\end{itemize}
Code, pre-registrations and all results, including failed tests, are released.
